
\section{Introduction}

This chapter presents the conclusions of the study, translates the empirical findings into strategic recommendations for Airtel Uganda, acknowledges the study's limitations, and proposes directions for future research. Section 6.2 states the conclusions against each objective and research hypothesis. Section 6.3 presents three evidence-based strategic recommendations for Airtel Uganda. Section 6.4 acknowledges the study's limitations. Section 6.5 proposes future research directions. Section 6.6 offers final remarks on the study's contribution.


\section{Conclusions}


\subsection{Conclusion on Specific Objective 1: Key Customer Attributes}

The first specific objective was to identify the key customer behavioural attributes and patterns most predictive of voice bundle purchase decisions among Airtel Uganda prepaid subscribers. This objective has been comprehensively achieved.

Applied to 1,458,900 real Airtel Uganda subscriber records over the period February to April 2026, Spearman Rank Correlation identified 34 of 35 features as significantly associated with value segment membership ($p < 0.05$), with IS\_SMARTPHONE producing the strongest correlation (r = -0.068) among non-revenue features, followed by BUNDLE\_SPEND\_30D (r = -0.067) and ACTIVE\_DAYS (r = -0.059). One-Way ANOVA confirmed that all 35 features exhibit significantly different mean values across the eight value segment groups, with VOICE\_INTENSITY (F = 27,892.34), IS\_SMARTPHONE (F = 25,891.16), BUNDLE\_SPEND\_30D (F = 23,594.82), and DISTINCT\_BUNDLES\_TRIED (F = 12,106.23) among the strongest non-revenue discriminators. Chi-Square testing confirmed that all 7 categorical features --- including HANDSET\_TYPE ($\chi^2$ = 181,566.52), MOST\_FREQ\_CHANNEL ($\chi^2$ = 59,471.19), and REGION ($\chi^2$ = 41,482.59) --- are significantly associated with segment. The one non-significant feature across all tests was UPGRADE\_RATE (Spearman r = 0.0004, p = 0.65), a novel finding indicating that upgrade frequency has no monotonic relationship with segment.

The central conclusion of Objective 1 is that smartphone ownership, recent bundle spend, voice call intensity, and bundle variety are the dominant behavioural signals for predicting voice bundle segment membership in the Ugandan prepaid market. These signals are observable in real time from Airtel's network infrastructure and do not require manual customer surveys or external data sources. H1	\textsubscript{1} is accepted: multiple customer behavioural attributes are significant predictors of voice bundle purchase decisions.


\subsection{Conclusion on Specific Objective 2: Model Selection and Performance}

The second specific objective was to test and evaluate suitable machine learning models for accurately predicting voice bundle purchase decisions, including appropriate data preprocessing and feature engineering. This objective has been fully achieved.

Four models were systematically evaluated on a hold-out test set of 446,599 records: Logistic Regression (F1 = 0.4465), Decision Tree (F1 = 0.6344), Random Forest (F1 = 0.6522), and XGBoost (F1 = 0.6611). XGBoost was identified as the best-performing model. The F1-Score range of 0.2146 across models confirms a statistically meaningful performance hierarchy. The XGBoost model's weighted accuracy of 68.0 percent represents a 5.4-fold improvement over a random baseline of 12.5 percent in the 8-class classification task, and its weighted precision of 0.670 means that segment-appropriate bundle recommendations are generated approximately 67 percent of the time --- compared to the current generic conversion rate of approximately 3 percent. H2	\textsubscript{1} is accepted: XGBoost significantly and meaningfully outperforms all other models.

The data preprocessing pipeline encompassing median imputation, capped SMOTE class balancing (limiting synthetic generation to a maximum of 50,000 samples per minority class), and the engineering of SPEND\_PER\_MINUTE and VOICE\_INTENSITY is concluded to be as critical a determinant of model performance as the choice of algorithm. SPEND\_PER\_MINUTE alone contributed 6.57 percent of the Random Forest model's predictive power --- more than CALLS (5.97\%) or MOBILE\_MONEY (4.64\%), despite being derived entirely from existing variables, confirming that thoughtful feature engineering from domain knowledge about prepaid behavioural patterns is a high-impact intervention (Singh et al., 2025; Yuming et al., 2024).


\subsection{Conclusion on Specific Objective 3: Business Impact and Strategic Value}

The third specific objective was to formulate strategic, evidence-based recommendations for Airtel Uganda to grow its business and improve customer satisfaction. The projected business impact analysis demonstrated a four- to five-fold improvement in bundle offer conversion rate (from approximately 3 percent to 12--15 percent) and a projected 9 percent improvement in Average Revenue Per User, grounded in the XGBoost model's weighted precision of 0.670 and consistent with documented outcomes from comparable global telecom recommendation deployments (McKinsey \& Company, 2022; 2024). H3	\textsubscript{1} is preliminarily accepted, subject to confirmation through a live, randomised A/B test deployment.

Taken together, the conclusions of this study establish that a feature-rich, behaviourally grounded machine learning recommendation system is both technically feasible and commercially compelling for Airtel Uganda. The study contributes the first published, empirically validated ML recommendation framework built on real transactional data from a Ugandan telecommunications operator, covering 1,458,900 subscribers across 738 distinct voice and combo bundle products, addressing a documented gap identified in Chapter Two.

In summary, the three conclusions of this study are stated directly. First,
customer behaviour predicts voice bundle segment: smartphone ownership, recent
bundle spend, voice call intensity, and bundle variety are the dominant signals,
all observable in real time from Airtel's network infrastructure
(H1\textsubscript{1} accepted). Second, XGBoost is the best model for this
classification task, achieving a 5.4-fold improvement over a random baseline on
an 8-class, 1.46-million-subscriber problem (H2\textsubscript{1} accepted). Third,
the model's outputs project commercially meaningful improvements in conversion rate
and ARPU, grounded in a weighted precision of 0.670 and consistent with global
telecom benchmarks, pending live A/B validation
(H3\textsubscript{1} preliminarily accepted). Together, these conclusions establish
that the primary barrier to personalised voice bundle recommendation at Airtel
Uganda is not data availability, not model performance, and not infrastructure
readiness --- it is the absence of the governance framework and live validation
experiment specified in the recommendations that follow.

\section{Strategic Recommendations for Airtel Uganda}

The following three strategic recommendations are grounded directly in the empirical findings of this study and designed to be operationally implementable within Airtel Uganda's existing technical infrastructure. The complete analytical codebase, datasets, and model outputs supporting these conclusions are documented in Appendix VI.


\subsection{Recommendation 1: Deploy a VOICE\_INTENSITY-Triggered Personalised Recommendation Engine}

The most actionable finding of this study is that VOICE\_INTENSITY --- calls per active day --- contributed 5.87 percent of Random Forest predictive power and produced the second-highest ANOVA F-statistic among non-revenue engineered features (F = 27,892.34). This has a direct and immediate operational implication: Airtel Uganda's recommendation delivery mechanism should be restructured from periodic, batch-run ARPU-banded campaigns to a dynamic, event-driven architecture that activates when a subscriber's call intensity crosses a defined behavioural threshold.

Specifically, it is recommended that Airtel Uganda deploy the XGBoost classification model as an automated real-time recommendation engine, configured to trigger a personalised bundle recommendation --- delivered via USSD push notification or SMS --- when a subscriber's VOICE\_INTENSITY (calls per active day) exceeds their 30-day rolling average by more than 20 percent. This indicates that the subscriber's current bundle is being consumed faster than their typical pattern, signalling active purchase intent. For subscribers whose SPEND\_PER\_MINUTE --- the second most important feature at 6.57 percent --- is above the median for their current value segment, the engine should additionally surface a higher-tier bundle recommendation, as the subscriber's efficiency profile suggests capacity for a more expensive product.

The projected commercial impact of this intervention --- a four- to five-fold improvement in bundle offer conversion rate compared to the current 3 percent baseline --- is grounded in the XGBoost model's weighted precision of 0.670 and is consistent with the broader literature on trigger-based marketing in telecoms (McKinsey \& Company, 2022; Wassouf et al., 2020). This approach directly operationalises the Push-Pull-Mooring logic of Chapter Two (Al-Mashraie et al., 2020): the VOICE\_INTENSITY signal identifies the precise moment at which the customer's bundle relationship is in flux, and the personalised recommendation serves as the mooring mechanism.


\subsection{Recommendation 2: Enrich Segmentation with Device Type and Channel Features}

The feature importance and statistical test results confirm that IS\_SMARTPHONE (Spearman r = -0.068; ANOVA F = 25,891.16; $\chi^2$ = 181,566.52) and \\MOST\_FREQ\_CHANNEL ($\chi^2$ = 59,471.19) are among the strongest non-revenue discriminators of bundle segment membership. Airtel Uganda's current ARPU-banded segmentation does not incorporate device generation or preferred purchase channel as segmentation variables --- a gap this study directly addresses.

It is recommended that Airtel Uganda enrich its segmentation framework by incorporating three additional real-time signals alongside VOICE\_INTENSITY and \\SPEND\_PER\_MINUTE: the subscriber's handset generation (2G, 3G, 4G, or 5G --- already available from the device master table); the subscriber's preferred transaction channel (USSD, App, or agent --- available from OCS transaction logs); and the subscriber's geographic region --- confirmed significant at $\chi^2$ = 41,482.59. A phased rollout is recommended, prioritising subscribers in the Kampala, Central 1, and Central 2 regions --- the geographic scope of this study --- as the initial deployment cohort, before national expansion.

This enriched segmentation framework does not require new data collection infrastructure. All three variables are already captured in Airtel Uganda's existing database systems (subscriber\_dimension, device\_master, ocsmonmaininternal, and geo\_master) and were used directly in this study's analysis. The incremental cost of incorporating them into the recommendation engine is therefore minimal relative to the projected commercial impact.

\subsection{Recommendation 3: Establish a Formal A/B Test and Model Governance Framework}
\label{sec:abt}

The projected business impact figures in this study are derived from a simulated pilot on the hold-out test set. Full confirmation of H3	\textsubscript{1} requires a live, randomised A/B test deployment. It is recommended that Airtel Uganda structure this as follows: within the priority subscriber cohort, randomly assign subscribers to a treatment group --- who receive personalised XGBoost model-driven recommendations at VOICE\_INTENSITY trigger points --- and a control group who continue to receive the standard undifferentiated product experience. Primary outcome metrics should include bundle purchase conversion rate, ARPU over the 90-day test period, and bundle repurchase frequency. The test should run for a minimum of 90 days to capture sufficient transaction volume for statistically significant results.

Alongside the A/B test, it is recommended that Airtel Uganda establish a formal Model Governance Framework comprising three components. First, a real-time performance monitoring dashboard tracking the XGBoost model's weighted F1-Score on a rolling weekly basis against the baseline of 0.6611 established in this study. Second, a quarterly scheduled retraining cycle in which the model is retrained on the most recent 90-day transaction window --- important because bundle portfolios evolve frequently (738 bundles were active in the February--April 2026 window) and subscriber behavioural patterns shift seasonally. Third, a closed-loop feedback mechanism logging each recommendation event and its outcome --- whether the subscriber purchased a recommended bundle, purchased a different bundle, or made no purchase --- to enable continuous improvement of the recommendation logic. This governance framework is consistent with best practice in production machine learning systems (Deldjoo et al., 2024; McKinsey \& Company, 2024).


\section{Limitations of the Study}

Several limitations are acknowledged. First, the classification task addressed an 8-class problem with severe class imbalance: the Silver segment comprised 53.3 percent of the dataset (776,970 subscribers) while the Diamond segment contained only 34 subscribers. Despite capped SMOTE oversampling, the minority classes --- particularly Diamond (F1 = 0.56), Ivory (F1 = 0.33), and Super Gold (F1 = 0.46) --- achieved lower per-class performance than the majority classes. The overall XGBoost F1-Score of 0.6611, while representing a 5.4-fold improvement over random, falls below the binary classification benchmarks ($\approx$ 0.82) commonly cited in the global literature --- a difference attributable to the task complexity rather than model inadequacy.

Second, the purchase timing dataset (PURCHASE\_TIMING\_FEB, MARCH, APRIL) contained 59,163,269 voice bundle transactions, of which the bundle ID filter matched 267 of 738 voice bundle IDs --- indicating that a significant portion of the purchase timing data contained bundle IDs not present in the current product catalogue, likely reflecting retired or promotional bundles. Future work should resolve this by using a historical rather than a current product catalogue for the join. Third, the dataset was confined to the Kampala, Central 1, and Central 2 regions, which may not fully capture the behavioural variation of subscribers in rural and peri-urban areas across Uganda. Fourth, the absence of a live A/B test deployment means that the projected business impact figures remain estimates grounded in model performance and external benchmarks, rather than directly observed commercial outcomes.


\section{Directions for Future Research}

The findings and limitations of this study suggest six concrete directions for
future research.

\textbf{Live A/B test validation.} The most immediate priority is the randomised
test specified in Recommendation 3 (Section~6.3.3), to be run for a minimum of
90 days in the Kampala, Central 1, and Central 2 regions. Primary outcome metrics
should include bundle purchase conversion rate, Average Revenue Per User, and
bundle repurchase frequency. This converts H3\textsubscript{1} from preliminary
support to empirical confirmation.

\textbf{Historical product catalogue integration.} Using a historical catalogue
snapshot covering the full February--April 2026 observation window would resolve
the bundle identifier mismatch --- 267 of 738 bundle identifiers matched the
current catalogue --- and unlock the full purchase history feature set, particularly
DISTINCT\_BUNDLES\_TRIED and BUNDLE\_SWITCH\_RATE, which could only be partially
computed in this study.

\textbf{Explainable AI (XAI) --- SHAP implementation.} SHAP (SHapley Additive
exPlanations) values \cite{ref_shap} should be incorporated into the production
recommendation pipeline to provide per-subscriber explanation of individual
predictions. This satisfies both internal commercial trust requirements and the
transparency provisions of Uganda's Data Protection and Privacy Act (2019)
\cite{ref47}.

\textbf{Cost-sensitive learning for minority classes.} The Diamond class comprised
only 34 real training instances and achieved a recall of 0.39 despite capped SMOTE.
Future work should explore class-weight penalties in XGBoost and Random Forest as
an alternative to oversampling for extreme minority classes, which would penalise
misclassification of rare but high-value subscribers more heavily during training.

\textbf{Deep Reinforcement Learning (DRL).} The XGBoost classifier is a static
model trained on historical patterns. A DRL agent would learn dynamically from
live customer responses, continuously updating its recommendation policy in real
time to maximise long-term customer lifetime value rather than predicting a single
segment label \cite{ref19}.

\textbf{Geographic expansion and seasonal modelling.} Extending the data scope to
Northern, Western, and Eastern Uganda would test whether the
VOICE\_INTENSITY and SPEND\_PER\_MINUTE trigger signals identified in this study
are nationally generalisable. Separately, extending the observation window to
12 months would enable the model to capture annual seasonal patterns --- December
festive spending, school-term income constraints, and agricultural income cycles
--- that the current 90-day window cannot represent.

For a researcher entering this space, the single highest-priority direction is the
live A/B test. The projected business impacts of this framework --- a conversion
rate improvement from 3\% to 12--15\% and a 9\% ARPU uplift (simulated pilot projection) --- are grounded in
model precision and external benchmarks, but they remain projections until a
randomised experiment confirms them empirically. Every subsequent research
direction --- DRL evolution, geographic expansion, XAI implementation, seasonal
modelling --- builds on a commercial baseline that only the A/B test can
establish. A researcher who confirms or refutes these projections through a live
deployment will make the most impactful single contribution to this body of work.

\section{Final Remarks}

This study set out to address a clearly identified gap in both academic literature and commercial practice: the absence of a published, empirically validated machine learning recommendation framework for the Ugandan prepaid voice bundle market. By analysing real transactional data from 1,458,900 Airtel Uganda subscribers across 738 voice and combo bundle products, applying a rigorous preprocessing and feature engineering pipeline, and systematically evaluating four machine learning models, the research has delivered a technically robust and commercially actionable framework that directly addresses this gap.

\subsection*{Contributions of this Study}

The study makes four distinct contributions. \textbf{Empirically}, it is the first
published ML recommendation framework validated on real, proprietary transactional
data from a Ugandan mobile network operator --- 1,458,900 subscribers across 738
products --- closing a gap that prior studies, which rely on public or synthetic
data, could not address. \textbf{Methodologically}, it demonstrates that a
classification-based recommendation paradigm is viable and deployable in
sparse-interaction prepaid telecom environments, and provides a reproducible
35-feature engineering pipeline and capped SMOTE protocol as a template for
similar studies. \textbf{Theoretically}, it establishes SPEND\_PER\_MINUTE ---
spending efficiency relative to usage --- as a stronger predictor of bundle tier
preference than absolute revenue, directly challenging the theoretical basis of
ARPU banding as a segmentation approach. \textbf{Practically}, it delivers a
deployment-ready roadmap --- a VOICE\_INTENSITY-triggered recommendation engine,
enriched segmentation variables, and a Model Governance Framework --- that requires
no new data collection infrastructure, enabling Airtel Uganda to act on these
findings without capital investment in new systems.

The study's central empirical contribution is the identification of SPEND\_PER\_MINUTE --- the ratio of total revenue to outgoing minutes consumed --- as the second most predictive feature overall (6.57\%), ahead of raw call counts and mobile money activity. This finding is not available in any prior study of Ugandan telecom data and directly challenges the adequacy of ARPU-banded segmentation: it is not how much a subscriber spends in absolute terms that most powerfully discriminates their bundle tier, but how efficiently they spend relative to their usage. Combined with VOICE\_INTENSITY (calls per active day; 5.87\%), these two engineered features provide Airtel Uganda with real-time, network-observable trigger signals for personalised recommendations that require no new data collection infrastructure.

The XGBoost model's F1-Score of 0.6611 on an 8-class, 1.46-million-subscriber classification task --- representing a 5.4-fold improvement over random and a 22-fold improvement over the current generic conversion rate- establishes a strong proof of concept for the ML-based personalisation approach in the Ugandan prepaid context. The full analytical codebase, datasets, and model outputs for this study are publicly available at https://github.com/RemmyBisimbeko/Data-Science/tree/main/SEM\%204/Thesis, enabling replication, extension, and adaptation by future researchers and practitioners working in related contexts across Sub-Saharan Africa and beyond.
